\documentclass[11pt]{article}
\usepackage{kristiania-workbook}

\begin{document}

\begin{center}
  {\LARGE\bfseries Set up your team, your lab \& your Microsoft Learn profile}\\[3pt]
  {\large Session 1 -- Introduction to Cloud Security}\\[5pt]
  {\small Exercise 1 \textperiodcentered\ Session 1 lab \textperiodcentered\ Estimated time: 60--90 min}
\end{center}
\vspace{2pt}\hrule\vspace{1.1em}

Three jobs today: (1) form the
\textbf{team you will keep for the whole semester and the exam}, (2) build the local Ubuntu laboratory
that is the red thread of this course, and (3) start the Microsoft Learn track that runs alongside it.

\par\medskip\noindent
\textbf{How the labs work (applies all semester).} Work through each lab \emph{in your group} and
discuss the answers together. At the next session you briefly talk the teacher through your answers --
\emph{verbally}, just so progress can be tracked (nothing is handed in, signed or graded here). The
review runs \emph{one session behind}, so you always have the week in between: Session 2 $\rightarrow$
Exercise 1; Session 3 $\rightarrow$ Exercise 2; \dots\ Session $N$ $\rightarrow$ Exercise $N\!-\!1$.
\par\medskip

\wbsection{1\quad Form your team (this is also your exam group)}
The home exam is done in groups, so the team you form now is the team you will be assessed with.

\note{Form a group of \textbf{3 to 5 people}. Fewer than 3 makes the workload heavy; more than 5 makes
coordination hard. You do \emph{not} register the members here -- what matters for the labs is your
group's \emph{progress}, not who is in it. The same group submits the final report, evidence appendix
and screen recording.}

\task{Pick a short \textbf{group name/ID}} (e.g.\ \cmd{OSL-07}) and use it on every answer document.

\note{\textbf{The answer documents (read once, apply all term).} Each week you fill in that session's
\textbf{answer document} --- a Word file (\cmd{SKY2100\_ExerciseN\_Answers.docx}) that mirrors this sheet:
every task, table and question here has a box there. Nothing else is collected. Keep one answer
document per exercise in your group's shared folder. The final home-exam report is built from these
weekly documents, so filling them in properly as you go \emph{is} the exam preparation.}

\wbsection{2\quad Build the local laboratory}

\task{Confirm virtualisation is enabled.} In your OS task manager (Windows: \cmd{Ctrl+Shift+Esc}
$\rightarrow$ Performance $\rightarrow$ CPU), check that \emph{Virtualization} is \textbf{Enabled};
otherwise enable Intel VT-x / AMD-V in BIOS/UEFI.

\task{Create and boot an Ubuntu VM} in VMware/VirtualBox ($\geq$2 vCPU, $\geq$4\,GB RAM,
$\geq$25\,GB disk). Open a terminal and run \cmd{whoami}, \cmd{uname -a}, \cmd{ip a}.

\caution{If the VM will not boot, do not get stuck: note it in the answer document, finish Section~3,
and ask an SA. You must not fall behind on the lab thread.}

\wbsection{3\quad Cloud transfer: Microsoft Learn}
We use the free \emph{Microsoft Learn} ``Introduction to Cloud Infrastructure'' series as the Azure
side of the course. (Module names verified August 2026.)

\task{Sign in and start Part 1 -- ``Describe cloud concepts''.}
\par\noindent\mbox{\href{https://learn.microsoft.com/training/paths/microsoft-azure-fundamentals-describe-cloud-concepts/}{learn.microsoft.com/training/paths/microsoft-azure-fundamentals-describe-cloud-concepts}}\\
Work through the first units: cloud computing, deployment models, and the shared-responsibility model.
Note in the answer document which units you completed.

\wbsection{4\quad Azure reflection}
Answer in your own words -- connect the training to today's lecture.

\task{In the shared-responsibility model, name one thing that is the \emph{provider's} job in SaaS but
\emph{your} job in IaaS.}

\task{The Microsoft training and the lecture both define ``the cloud''. Give the \emph{one} difference
you found most useful between how a vendor (Microsoft) frames it and how the course frames it (from a
security view).}

\aha{The two halves of today differed in who was responsible. You build and maintain the local VM
yourself. On Azure you look after only your account and your data, while the provider runs the hardware,
the hypervisor and the platform. That division of labour is the shared-responsibility model, and how
much is your job depends on the service model: IaaS, PaaS or SaaS.}

\wbsection{5\quad Knowledge check}
Write short answers in the space provided. Tick the correct box for multiple-choice items.

\begin{enumerate}
\item Which is \textbf{not} one of the five essential cloud characteristics?\\
      \cbox on-demand self-service \quad \cbox broad network access \quad
      \cbox unlimited free storage \quad \cbox rapid elasticity
\hint{Four are genuine NIST characteristics; one promises something no paid service delivers.}
\item Name the three CIA properties and give a one-line meaning for each.
\hint{One keeps secrets, one prevents undetected tampering, one keeps the service reachable.}
\item Explain the difference between a \textbf{threat}, a \textbf{risk} and a \textbf{control}, with
      one example of your own.
\hint{One is the potential cause, one is likelihood times impact, one is the measure that lowers it.}
\item Which service model gives the customer the \emph{most} control of the stack?\\
      \cbox SaaS \quad \cbox PaaS \quad \cbox IaaS \quad \cbox SecaaS
\hint{The model where you manage the operating system and everything above it.}
\item Match each to on-prem / IaaS / PaaS / SaaS: ``I manage only my data and users'' \,=\,
      \rule{2.2cm}{0.3pt}; \quad ``I manage the OS upward'' \,=\, \rule{2.2cm}{0.3pt}.
\hint{The least you manage is 'only data and users'; managing the OS upward is the other extreme.}
\item In one sentence: why is securing the cloud different from securing an internal server?
\hint{The provider now shares the stack with you — think trust boundaries and shared hardware.}
\item Define \textbf{Security as a Service (SecaaS)} and give one example.
\hint{Security delivered like any other cloud service: a managed WAF, SIEM or DDoS protection.}
\item Name the four deployment models and say which one dominates in practice.
\hint{Public, private, community — and the blend that dominates in practice.}
\item The group home exam rewards primarily\dots\\
      \cbox presentation polish \quad \cbox consistency between report and evidence \quad
      \cbox number of screenshots \quad \cbox report length
\hint{It is not about polish; the report and the evidence have to agree.}
\item True or false: \emph{undisclosed} AI use is the problem, not AI use itself. Briefly justify.
\hint{The declaration targets the undisclosed kind; an LLM cannot produce your lab evidence.}
\item Pick one CIA property and name a \emph{later session} where it is the main focus.
\hint{Confidentiality → the TLS and access sessions; Integrity → logs and timeline; Availability → load or backup.}
\item (Reflection) Write one question about cloud security you want answered by the end of the course.
\end{enumerate}

\signoff

\clearpage
\solheader
{\LARGE\bfseries\color{kred} Solutions}\\[2pt]
{\small Work through the lab first, then check here. Model answers are indicative — equivalent reasoning is fine.}\par\vspace{8pt}
\noindent\textbf{Note:} Model answers are indicative -- accept equivalent reasoning. Multiple-choice
answers are in \textbf{bold}.

\wbsection{Knowledge check}
\begin{enumerate}
\item \textbf{Unlimited free storage} is not a characteristic. The five are: on-demand self-service,
      broad network access, resource pooling, rapid elasticity, measured service.
\item \textbf{Confidentiality} -- disclosed only to the authorised; \textbf{Integrity} -- not altered
      without authorisation/detection; \textbf{Availability} -- reachable when needed.
\item \textbf{Threat} = a potential cause of an unwanted incident; \textbf{Risk} = likelihood
      $\times$ impact; \textbf{Control} = a measure that reduces the risk. E.g.\ threat = request
      flood; risk = high for a public service; control = rate limiting / WAF.
\item \textbf{IaaS} gives the customer the most control of the stack.
\item ``Only my data and users'' = \textbf{SaaS}; ``OS upward'' = \textbf{IaaS}.
\item Because \emph{control is shared} with the provider: trust boundaries are logical not physical,
      you share infrastructure (multi-tenancy), and depend on an internet-facing control plane.
\item \textbf{SecaaS} = security capabilities consumed as a cloud service; e.g.\ a managed WAF, cloud
      monitoring/SIEM, or managed DDoS protection.
\item \textbf{Public, Private, Community, Hybrid}; \textbf{Hybrid} dominates in practice.
\item \textbf{Consistency between report and evidence.}
\item \textbf{True} -- the AI declaration penalises \emph{undisclosed} use; an LLM cannot produce your
      lab evidence.
\item Any valid pair, e.g.\ Confidentiality $\rightarrow$ S6--7; Integrity $\rightarrow$ S9--10;
      Availability $\rightarrow$ S3/S11.
\item Open reflection -- no fixed answer.
\end{enumerate}

\wbsection{References}
CSA Security Guidance v5, \textbf{Domain~1: Cloud Computing Concepts \& Architectures} (verified);
Comer Ch.~1--3 (verified); J\o sang Ch.~1 (CIA, risk -- syllabus, page numbers not verified);
Lecture~1.

\end{document}
